
\documentclass{article}
\usepackage{colortbl}
\usepackage{makecell}
\usepackage{multirow}
\usepackage{supertabular}

\begin{document}

\newcounter{utterance}

\centering \large Interaction Transcript for game `ifeval', experiment `ifeval', episode 176 with qwen3.5{-}27b.
\vspace{24pt}

{ \footnotesize  \setcounter{utterance}{1}
\setlength{\tabcolsep}{0pt}
\begin{supertabular}{c@{$\;$}|p{.15\linewidth}@{}p{.15\linewidth}p{.15\linewidth}p{.15\linewidth}p{.15\linewidth}p{.15\linewidth}}
    \# & \multicolumn{2}{c}{Player} && \multicolumn{2}{c}{Game Master} \\
    \hline

    \theutterance \stepcounter{utterance}  
    & & & \multicolumn{4}{p{0.6\linewidth}}{
        \cellcolor[rgb]{0.9,0.9,0.9}{
            \makecell[{{p{\linewidth}}}]{
                \texttt{\tiny{[P1$\langle$GM]}}
                \texttt{I'm a new puppy owner and I'm looking for some advice on how to train my puppy. Can you help me? Give me a few options. In particular, I need you to end your response with "Which one you choose?".} \\
            }
        }
    }
     \\ \\

    \theutterance \stepcounter{utterance}  
    & \multicolumn{4}{p{0.6\linewidth}}{
        \cellcolor[rgb]{0.561,0.737,0.561}{
            \makecell[{{p{\linewidth}}}]{
                \texttt{\tiny{[P1$\rangle$GM]}}
                \texttt{Congratulations on your new puppy! Training is a wonderful journey that builds a strong bond between you and your dog. Here are three effective approaches you can consider:} \\
\\ 
\texttt{1.  **Positive Reinforcement**: This is the most widely recommended method by modern trainers. It involves rewarding your puppy with treats, praise, or play immediately after they perform a desired behavior. This encourages them to repeat the good actions because they associate them with something pleasant.} \\
\texttt{2.  **Clicker Training**: A specific form of positive reinforcement that uses a small device called a "clicker" to mark the exact moment your puppy does something right. The click sound acts as a precise signal that a reward is coming, which can help your puppy understand exactly which behavior earned the treat.} \\
\texttt{3.  **Lifestyle Integration (Routine{-}Based)**: Instead of formal "training sessions," this approach focuses on embedding commands and boundaries into your daily life. You practice "sit" before meals, "stay" before opening the door, and "come" during playtime, making training a natural part of your day rather than a separate chore.} \\
\\ 
\texttt{Each method has its strengths, and many owners find success by combining elements of them. Which one you choose?} \\
            }
        }
    }
    & & \\ \\

    \theutterance \stepcounter{utterance}  
    & & & \multicolumn{2}{p{0.3\linewidth}}{
        \cellcolor[rgb]{0.9,0.9,0.9}{
            \makecell[{{p{\linewidth}}}]{
                \texttt{\tiny{[GM$|$GM]}}
                {'startend:end_checker': {'end_phrase': 'Which one you choose?'}}
            }
        }
    }
    & & \\ \\

    \theutterance \stepcounter{utterance}  
    & & & \multicolumn{2}{p{0.3\linewidth}}{
        \cellcolor[rgb]{0.9,0.9,0.9}{
            \makecell[{{p{\linewidth}}}]{
                \texttt{\tiny{[GM$|$GM]}}
                \texttt{game\_result = WIN} \\
            }
        }
    }
    & & \\ \\

\end{supertabular}
}

\end{document}
